\documentclass[aspectratio=169]{beamer}

%================================================
% PAQUETES
%================================================

\usepackage[utf8]{inputenc}
\usepackage[T1]{fontenc}
\usepackage[spanish,es-noshorthands]{babel}
\usepackage{lmodern}
\usepackage{amsmath,amssymb,amsfonts}
\usepackage{booktabs}
\usepackage{array}
\usepackage{multirow}
\usepackage{enumitem}
\usepackage{tikz}
\usepackage{tcolorbox}
\usepackage{graphicx}
\usepackage{multicol}
\usepackage{siunitx}
\usepackage{ragged2e}

\sisetup{
	output-decimal-marker={,}
}

%================================================
% COLORES UST MODERNO
%================================================

\definecolor{Primary}{HTML}{0B6E4F}
\definecolor{Secondary}{HTML}{1B4332}
\definecolor{Accent}{HTML}{2D6A4F}
\definecolor{Soft}{HTML}{F4F7F5}
\definecolor{Dark}{HTML}{1C1C1C}

\setbeamercolor{title}{fg=white,bg=Primary}
\setbeamercolor{frametitle}{fg=white,bg=Primary}
\setbeamercolor{structure}{fg=Primary}
\setbeamercolor{block title}{fg=white,bg=Primary}
\setbeamercolor{block body}{bg=Soft}
\setbeamercolor{normal text}{fg=Dark}

\setbeamertemplate{navigation symbols}{}
\setbeamertemplate{blocks}[rounded][shadow=false]

%================================================
% FOOTLINE
%================================================

\setbeamertemplate{footline}{
	\leavevmode
	\hbox{
		\begin{beamercolorbox}[wd=.82\paperwidth,ht=2.5ex,dp=1ex,leftskip=.3cm]{author in head/foot}
			\color{white}\small Fundamentos de Matemática Básica
		\end{beamercolorbox}
		\begin{beamercolorbox}[wd=.18\paperwidth,ht=2.5ex,dp=1ex,center]{date in head/foot}
			\color{white}\small \insertframenumber/\inserttotalframenumber
		\end{beamercolorbox}
	}
}

\setbeamercolor{author in head/foot}{bg=Secondary}
\setbeamercolor{date in head/foot}{bg=Primary}

%================================================
% DATOS
%================================================

\title{Clase Aplicación de ecuaciones}
\subtitle{Ecuaciones de Primer Grado Aplicadas a la Agronomía}
\author{Prof. Luis Tulio González Alcaino}
\institute{Universidad Santo Tomás\\Departamento de Ciencias Básicas}
\date{2026}

%================================================
% DOCUMENTO
%================================================

\begin{document}
	
	%================================================
	\begin{frame}
		\titlepage
	\end{frame}
	
	%================================================
	\begin{frame}{Objetivo de la clase}
		\begin{block}{Aprendizaje esperado}
			Aplicar ecuaciones de primer grado para modelar y resolver
			situaciones contextualizadas en Agronomía e Ingeniería Agrícola,
			distinguiendo ecuaciones enteras, fraccionarias y literales,
			analizando patrones, relaciones e inconsistencias en los resultados.
		\end{block}
	\end{frame}
	
	%================================================
	\begin{frame}{Ruta de la clase}
		\tableofcontents
	\end{frame}
	
	%================================================
	\section{Modelación Matemática}
	
	\begin{frame}{¿Por qué modelamos?}
		
		En Agronomía, muchas decisiones técnicas se toman a partir de:
		
		\begin{itemize}
			\item dosis de fertilización
			\item consumo hídrico
			\item costos de producción
			\item rendimiento por hectárea
			\item eficiencia de riego
			\item control fitosanitario
		\end{itemize}
		
		Todo esto puede representarse mediante ecuaciones.
		
		\begin{block}{Idea clave}
			Resolver una ecuación significa encontrar el valor desconocido
			que hace verdadera una situación real.
		\end{block}
		
	\end{frame}
	
	%================================================
	\section{Ecuaciones Enteras}
	
	\begin{frame}{Ecuaciones Enteras}
		
		Forma general:
		
		\[
		ax+b=c
		\]
		
		donde:
		
		\begin{itemize}
			\item \(a\): coeficiente
			\item \(x\): incógnita
			\item \(b\): término independiente
			\item \(c\): valor conocido
		\end{itemize}
		
	\end{frame}
	
	%================================================
	\begin{frame}{Ejemplo 1: Fertilización nitrogenada}
		
		Un cultivo requiere 240 kg de nitrógeno.
		
		Ya se aplicaron 60 kg.
		
		Cada saco aporta 15 kg.
		
		¿Cuántos sacos deben agregarse?
		
		\[
		15x + 60 = 240
		\]
		
	\end{frame}
	
	%================================================
	\begin{frame}{Resolución paso a paso}
		
		\[
		15x + 60 = 240
		\]
		
		\pause
		
		\[
		15x = 240 - 60
		\]
		
		\pause
		
		\[
		15x = 180
		\]
		
		\pause
		
		\[
		x = \frac{180}{15}
		\]
		
		\pause
		
		\[
		x = 12
		\]
		
		\begin{block}{Interpretación}
			Se requieren 12 sacos de fertilizante.
		\end{block}
		
	\end{frame}
	
	%================================================
	\section{Ecuaciones Fraccionarias}
	
	\begin{frame}{Ecuaciones Fraccionarias}
		
		Contienen fracciones algebraicas o numéricas.
		
		Ejemplo general:
		
		\[
		\frac{x}{a} + b = c
		\]
		
		Se recomienda eliminar denominadores primero.
		
	\end{frame}
	
	%================================================
	\begin{frame}{Ejemplo 2: Sistema de riego}
		
		Cada planta necesita 2,5 litros diarios.
		
		Un estanque entrega en total 450 litros.
		
		Si ya se consumieron 75 litros:
		
		\[
		\frac{x}{2,5} + 75 = 255
		\]
		
		Determinar la cantidad de plantas abastecidas.
		
	\end{frame}
	
	%================================================
	\begin{frame}{Resolución}
		
		\[
		\frac{x}{2,5} + 75 = 255
		\]
		
		\pause
		
		\[
		\frac{x}{2,5} = 180
		\]
		
		\pause
		
		\[
		x = 180(2,5)
		\]
		
		\pause
		
		\[
		x = 450
		\]
		
		\begin{block}{Conclusión}
			El sistema abastece correctamente a 450 plantas.
		\end{block}
		
	\end{frame}
	
	%================================================
	\section{Ecuaciones Literales}
	
	\begin{frame}{Ecuaciones Literales}
		
		Son ecuaciones donde además de la incógnita aparecen letras
		como parámetros.
		
		Ejemplo:
		
		\[
		A = bx + c
		\]
		
		Se despeja una variable respecto de otra.
		
	\end{frame}
	
	%================================================
	\begin{frame}{Ejemplo 3: Producción agrícola}
		
		La producción total está dada por:
		
		\[
		P = rx + 120
		\]
		
		donde:
		
		\begin{itemize}
			\item \(P\): producción total
			\item \(r\): rendimiento por hectárea
			\item \(x\): hectáreas cultivadas
		\end{itemize}
		
		Despejar \(x\).
		
	\end{frame}
	
	%================================================
	\begin{frame}{Resolución}
		
		\[
		P = rx + 120
		\]
		
		\pause
		
		\[
		P - 120 = rx
		\]
		
		\pause
		
		\[
		x = \frac{P-120}{r}
		\]
		
		\begin{block}{Interpretación}
			Permite estimar superficie cultivada
			a partir de producción observada.
		\end{block}
		
	\end{frame}
	
	%================================================
	\section{Patrones e Inconsistencias}
	
	\begin{frame}{Detectar inconsistencias}
		
		No basta resolver.
		
		Debemos validar:
		
		\begin{itemize}
			\item unidad correcta
			\item signo correcto
			\item magnitud razonable
			\item coherencia agronómica
		\end{itemize}
		
		\begin{block}{Ejemplo}
			Si obtenemos:
			
			\[
			x=-15
			\]
			
			y \(x\) representa número de sacos,
			el resultado es inconsistente.
		\end{block}
		
	\end{frame}
	
	%================================================
	\begin{frame}{Análisis de patrones}
		
		Observe:
		
		\[
		10x+100=200
		\]
		
		\[
		20x+100=300
		\]
		
		\[
		30x+100=400
		\]
		
		\pause
		
		¿Qué patrón aparece?
		
		\pause
		
		\[
		x=10
		\]
		
		en todos los casos.
		
		\begin{block}{Conclusión}
			Existe regularidad estructural,
			no solo cálculo aislado.
		\end{block}
		
	\end{frame}
	
	%================================================
	\section{Ejercicios Propuestos}
	
	\begin{frame}{Ejercicio 1}
		
		Una parcela necesita 350 kg de fertilizante.
		
		Ya existen 110 kg almacenados.
		
		Cada saco contiene 20 kg.
		
		Plantee y resuelva la ecuación.
		
	\end{frame}
	
	%================================================
	\begin{frame}{Ejercicio 2}
		
		Una bomba de riego distribuye agua de forma uniforme.
		
		Cada sector consume:
		
		\[
		\frac{x}{4} + 30 = 90
		\]
		
		Determine \(x\).
		
	\end{frame}
	
	%================================================
	\begin{frame}{Ejercicio 3}
		
		Despejar \(x\):
		
		\[
		C = px + 45
		\]
		
		Interpretar el resultado en contexto agrícola.
		
	\end{frame}
	
	%================================================
	\begin{frame}{Cierre}
		
		\begin{block}{Síntesis}
			Hoy aprendimos a:
			
			\begin{itemize}
				\item modelar situaciones reales
				\item resolver ecuaciones enteras
				\item resolver ecuaciones fraccionarias
				\item trabajar ecuaciones literales
				\item detectar inconsistencias y patrones
			\end{itemize}
			
		\end{block}
		
		\begin{center}
			\Large
			La matemática no solo resuelve:\\
			permite tomar mejores decisiones técnicas.
		\end{center}
		
	\end{frame}
	
	%================================================
	
\end{document}